\clearpage
\section{Especificação formal M7: idempotência}
\label{sec:formal-m7}

M7 define a identidade lógica de uma operação por
\texttt{(uid, tipo\_operacao, payload\_hash)}. O hash é SHA-256 de
\texttt{tipo\_operacao + "\\n" + canonicalize(payload)}; objetos são ordenados
recursivamente, arrays preservam ordem, e \texttt{null} permanece distinto de
campo ausente. O \texttt{idOperacao} é reutilizado em retry da mesma intenção,
mas não entra no hash.

\subsection*{Responsabilidade das camadas}
CUE verifica o contrato estrutural e os estados de processamento. Rust fornece
a canonicalização executável, vetores de teste, hashes e validação do receipt.
Alloy verifica deduplicação do efeito observado, retries, eventos,
materializações substitutivas e transições coerentes.

\begingroup\small\raggedright
% Gerado por lcqui-spec-doc; não editar.
\subsection*{M7\_Operacao\_Idempotente --- formalização executável M7}
Identidade, estados e deduplicação estrutural; SHA-256/canonicalização são validados por Rust.
\begin{description}
\item[id\_operacao] Tipo: string. Obrigatório: sim. Aceita nulo: não.
SQL: VARCHAR(100).
Máximo: 100 caracteres.
\item[uid] Tipo: string. Obrigatório: sim. Aceita nulo: não.
SQL: VARCHAR(100).
Máximo: 100 caracteres.
\item[tipo\_operacao] Tipo: string. Obrigatório: sim. Aceita nulo: não.
SQL: VARCHAR(100).
Máximo: 100 caracteres.
\item[payload\_hash] Tipo: string. Obrigatório: sim. Aceita nulo: não.
SQL: TEXT.
Máximo: 0 caracteres.
\item[status] Tipo: enum. Obrigatório: sim. Aceita nulo: não.
SQL: ENUM.
Valores: PENDENTE, CONCLUIDA, FALHOU.
Máximo: 0 caracteres.
\end{description}

% Gerado por lcqui-spec-doc; não editar.
\subsection*{Evidência formal M7 --- idempotência}
Alloy verifica identidade, retry, deduplicação e materialização substitutiva. Rust é a referência determinística para proveniência e canonicalização; isto não prova exactly-once da infraestrutura, TOCTOU ou deadlock.
\begin{description}
\item[M7-INV-001] IdentidadeUnicaNaoDuplicaEfeito: UNSAT.\newline Escopo: \texttt{check IdentidadeUnicaNaoDuplicaEfeito for 4}.
\item[M7-INV-002] RetryNaoReexecuta: UNSAT.\newline Escopo: \texttt{check RetryNaoReexecuta for 6}.
\item[M7-INV-003] ReusoIncompativelRejeitado: UNSAT.\newline Escopo: \texttt{check ReusoIncompativelRejeitado for 6}.
\item[M7-INV-004] EventoDeduplicado: UNSAT.\newline Escopo: \texttt{check EventoDeduplicado for 6}.
\item[M7-INV-005] MaterializacaoSubstitutiva: UNSAT.\newline Escopo: \texttt{check MaterializacaoSubstitutiva for 6}.
\item[M7-INV-006] TransicoesPreservamCoerencia: UNSAT.\newline Escopo: \texttt{check TransicoesPreservamCoerencia for 6}.
\item[M7-WIT-007] WitnessPrimeiraExecucao: SAT.\newline Escopo: \texttt{run WitnessPrimeiraExecucao for 4}.
\item[M7-WIT-008] WitnessRetry: SAT.\newline Escopo: \texttt{run WitnessRetry for 6 but exactly 3 Store}.
\item[M7-WIT-009] WitnessEventoDuplicado: SAT.\newline Escopo: \texttt{run WitnessEventoDuplicado for 6}.
\item[M7-WIT-010] WitnessMaterializacao: SAT.\newline Escopo: \texttt{run WitnessMaterializacao for 6}.
\end{description}
UNSAT para check indica ausência de contraexemplo no escopo declarado; SAT para run indica testemunha encontrada. A evidência não certifica a implementação Firebase/backend.

\par\endgroup

\subsection*{Limites}
Idempotência e deduplicação do efeito observado não equivalem a exactly-once da
infraestrutura. Este modelo não prova ausência de TOCTOU, deadlock, contenção ou
loops de triggers; esses temas permanecem fora de M7.
